\section{Recomendaciones}
\label{sec:recomendaciones}

\begin{table}[htbp]
    \centering
    \caption{Recomendaciones derivadas del an\'{a}lisis de la l\'{i}nea SIM-014.}
    \label{tab:recomendaciones}
    \setcounter{itemcount}{0}
    \begin{tabularx}{\textwidth}{@{}NlX@{}}
        \toprule
        \multicolumn{1}{c}{\textbf{Item}} & \textbf{Categoria} & \textbf{Recomendacion} \\
        \midrule
        & Construcci\'{o}n & Construir la soporter\'{i}a estrictamente conforme al isom\'{e}trico PCF112: la restricci\'{o}n lateral del nodo 40 y el conjunto de restricciones del nodo 105 (gu\'{i}as con holgura de 25~\si{mm} y soporte vertical +Y) son parte integral del comportamiento verificado; cualquier cambio de posici\'{o}n, tipo u holgura de soporte invalida el presente an\'{a}lisis. \\
        & Ingenier\'{i}a & Habilitar en el job el reporte de cargas de las conexiones de los extremos del tramo y remitirlas al responsable del TMS para confirmar la aceptaci\'{o}n de las cargas en su conexi\'{o}n de descarga. \\
        & Ingenier\'{i}a & Ante cualquier modificaci\'{o}n del modelo, verificar la asignaci\'{o}n uniforme del material ASTM A312 TP 304L en todos los componentes y la definici\'{o}n de tipo de los soportes, seg\'{u}n las advertencias de modelado declaradas en la secci\'{o}n~\ref{sec:alcance}. \\
        & Documentaci\'{o}n & Conservar el job CAESAR P2603-PR-SIM-014 y el PCF corregido como registro del an\'{a}lisis; ante cualquier modificaci\'{o}n de la l\'{i}nea (di\'{a}metros, recorrido, condiciones de operaci\'{o}n), re-ejecutar el modelo y re-emitir este informe con una nueva revisi\'{o}n. \\
        \bottomrule
    \end{tabularx}
\end{table}
